\documentclass[10pt,a4paper]{amsart}
\usepackage[margin=19mm]{geometry}
\usepackage{amsmath,amssymb,mathtools}
\usepackage[scr=rsfs,cal=euler]{mathalfa}
\usepackage{xfrac}
\usepackage[unicode,hidelinks,bookmarks=false]{hyperref}
\newcommand{\ie}{i.e.\ }
\newcommand{\cf}{cf.\ }
\newcommand{\resp}{respectively\ }
\newcommand{\Subsection}{\subsection}
\providecommand{\nobreakdash}{\nobreak\hbox{-}}
\setlength{\emergencystretch}{2em}
\multlinegap0pt
\usepackage{bm}
\usepackage{bbm}
\usepackage{dsfont}
\usepackage{xspace}
\usepackage{cancel}
\usepackage{mathtools}
\usepackage{cleveref}
\usepackage{amsthm}
\makeatletter
\@ifundefined{theorem}{\newtheorem{theorem}{Theorem}}{}
\@ifundefined{lemma}{\newtheorem{lemma}[theorem]{Lemma}}{}
\makeatother
\newcommand{\Q}{\mathbb Q}
\newcommand{\Z}{\mathbb Z}
\title[Reconstruction of Torsion-Free Abelian Groups from Rational ]{Reconstruction of Torsion-Free Abelian Groups from Rational Group Fields}
\author{Jinyu Lin, Xiaodong Wang}
\date{Source: arXiv:2608.10381v1 (CC BY 4.0)}
\begin{document}
\maketitle

\noindent\emph{Source and licence.} Reconstruction of Torsion-Free Abelian Groups from Rational Group Fields. Version arXiv:2608.10381v1 (CC BY 4.0), distributed under the Creative Commons Attribution 4.0 International licence, as recorded in the arXiv licence metadata for this exact version and shown on the abstract page https://arxiv.org/abs/2608.10381v1. Full rights notice: licenses/arxiv-2608.10381-CC-BY-4.0.txt.\par\smallskip

\noindent\textit{Editorial scope.} Counted unit: the Lemma whose source label is \texttt{lem:puiseux-free} in the article (source0.tex lines 309--315, source label \texttt{lem:puiseux-free}), delivered together with the article's own complete original proof of it (source0.tex lines 317--463, one \texttt{proof} environment of 147 lines). No mathematics is written, repaired or re-derived: statement and proof are the article's own text line for line. Required context retained uncounted: none. Every definition, display and label of those lines is retained unchanged. Omitted from this package and not redistributed: 1--308, 316--316, 464--995. Nothing inside a retained excerpt is omitted or altered: every retained body is delivered exactly as the article's lines are written, so no omission receipt of any kind accompanies this package. Citations: the retained excerpts carry no citation command at all, so no bibliography entry of the article is transcribed and no reference is redistributed. Reference adaptation: none was needed and none was made; every reference inside the delivered unit resolves inside it, so no unresolved reference or citation is produced. Printed slips kept literal: no printed word, symbol, hypothesis or number of the counted statement or of its proof was altered. Layout and class: the article's own class is \texttt{article} with packages including geometry, amsmath, amssymb, amsthm, mathtools, microtype, enumitem, xcolor, hyperref, cleveref; the wrapper is the lane's pinned \texttt{amsart} editorial wrapper at 10pt on A4, which restates the theorem-like environments the retained text needs and adds the article's own macro definitions that the retained text needs (\texttt{\textbackslash Q}, \texttt{\textbackslash Z}), with extra wrapper packages (bm, bbm, dsfont, xspace, cancel, mathtools, cleveref). The delivered numbering is the unit's own. Assets: no retained span contains a figure environment, a graphics-inclusion command, a PDF-page inclusion, a table or a TikZ picture; no image file is copied into this package, the package declares no asset dependency and no image file is redistributed with it. Licence: version arXiv:2608.10381v1 of this article is distributed under the Creative Commons Attribution 4.0 International licence, as recorded in the arXiv licence metadata for this exact version and shown on the abstract page https://arxiv.org/abs/2608.10381v1. A full rights notice is delivered at licenses/arxiv-2608.10381-CC-BY-4.0.txt. Characters: every retained excerpt is pure ASCII, so the delivered engine, XeLaTeX with its default Latin Modern text font, reports no missing character.
\begin{lemma}[Puiseux quotient]\label{lem:puiseux-free}
If $\operatorname{char}F=0$, then
\[
  P(F,t):=F(t^{\Q})^\times/(F^\times t^{\Q})
\]
is a free abelian group.
\end{lemma}

\begin{proof}
Put $z_n=t^{1/n!}$ and
\[
  P_n=F(z_n)^\times/(F^\times z_n^{\Z}).
\]
Let $\mathcal I_n$ be the set of monic irreducible polynomials in the
polynomial ring $F[z_n]$, with the polynomial $z_n$ omitted.  For each
$r\in\mathcal I_n$, let $v_r$ be the usual exponent valuation in the unique
factorization of a rational function.  Every $f\in F(z_n)^\times$ has a
unique expression
\begin{equation}\label{eq:rational-factorization}
  f=c z_n^m\prod_{r\in\mathcal I_n}r(z_n)^{a_r},
\end{equation}
where $c\in F^\times$, $m,a_r\in\Z$, and only finitely many $a_r$ are
nonzero.  Therefore the map
\[
  [f]\longmapsto (v_r(f))_{r\in\mathcal I_n}
\]
is a well-defined isomorphism
\begin{equation}\label{eq:Pn-free}
  P_n\xrightarrow{\ \sim\ }\Z^{(\mathcal I_n)}.
\end{equation}

We next determine the transition map under these bases.  Put $d=n+1$ and
write $Y=z_{n+1}$, so that
\begin{equation}\label{eq:stage-substitution}
  z_n=Y^d.
\end{equation}
Fix a monic irreducible $p(Z)\in F[Z]$ with $p\neq Z$.  Then $p(0)\neq0$.
Because $F$ has characteristic zero, $p$ is separable and
$\gcd(p,p')=1$.  Choose $A,B\in F[Z]$ satisfying
\[
  A(Z)p(Z)+B(Z)p'(Z)=1.
\]
After substituting $Z=Y^d$, this becomes
\begin{equation}\label{eq:bezout-substitution}
  A(Y^d)p(Y^d)+B(Y^d)p'(Y^d)=1.
\end{equation}
The formal derivative of $p(Y^d)$ is
\[
  \frac{d}{dY}p(Y^d)=dY^{d-1}p'(Y^d).
\]
Equation \cref{eq:bezout-substitution} shows that $p(Y^d)$ is coprime to
$p'(Y^d)$; its nonzero constant term shows that it is coprime to $Y$; and
$d\neq0$ in $F$.  Hence $p(Y^d)$ is coprime to its derivative and is
square-free.

If $p_1$ and $p_2$ are distinct monic irreducibles, then they are coprime.
Substitution into a Bezout identity for $p_1$ and $p_2$ shows that
$p_1(Y^d)$ and $p_2(Y^d)$ are coprime.  Thus no target irreducible can
occur above two distinct source irreducibles.

Conversely, let $q(Y)\neq Y$ be a monic irreducible polynomial over $F$,
and write $E=F[Y]/(q)$ and $\bar y=Y+(q)\in E$.  Since $q\neq Y$, we have
$\bar y\neq0$.  Let $p(Z)$ be the monic minimal polynomial of $\bar y^d$
over $F$.  Then $p$ is irreducible, $p\neq Z$, and
$p(\bar y^d)=0$.  The last equality is precisely the assertion that
$q(Y)$ divides $p(Y^d)$.  If $q$ also divides $\widetilde p(Y^d)$ for a
monic irreducible $\widetilde p$, then
$\widetilde p(\bar y^d)=0$, so minimality implies
$p=\widetilde p$.  Therefore each target irreducible lies above exactly one
source irreducible.

For $p\in\mathcal I_n$, let
\[
  S_p=\{q\in\mathcal I_{n+1}:q(Y)\text{ divides }p(Y^d)\}.
\]
Each $S_p$ is a nonempty finite set, the sets $S_p$ are pairwise disjoint,
and the preceding paragraph proves that they partition
$\mathcal I_{n+1}$.  Square-freeness gives
\[
  p(Y^d)=\prod_{q\in S_p}q(Y).
\]
Consequently, under \cref{eq:Pn-free}, the transition homomorphism sends
the basis vector $p$ to
\begin{equation}\label{eq:transition-vector}
  \sum_{q\in S_p}q.
\end{equation}
For each $p$, choose one element $q_p\in S_p$.  The family
\[
  \left\{\sum_{q\in S_p}q\right\}
  \cup\{q:q\in S_p\setminus\{q_p\}\}
\]
is a basis of $\Z^{(S_p)}$: relative to the original basis, replacing
$q_p$ by the sum has determinant $1$.  Thus
\[
  \Z^{(S_p)}=
  \Z\left(\sum_{q\in S_p}q\right)
  \oplus\Z^{(S_p\setminus\{q_p\})}.
\]
Taking the direct sum over all $p$ proves that
$P_n\to P_{n+1}$ is injective, that its image is a direct summand, and that
its cokernel is free abelian.

It remains to justify that these stage groups really form an increasing
union inside the desired quotient.  Let
\[
  \iota_n:P_n\longrightarrow
  P(F,t)=F(t^\Q)^\times/(F^\times t^\Q)
\]
be induced by the field inclusion.  Suppose $[f]\in\ker\iota_n$.  Then
$f=c t^q$ for some $c\in F^\times$ and $q\in\Q$.  Write
$f=a/b$ with nonzero $a,b\in F[z_n]$.  In the group algebra $F[\Q]$ we
then have
\begin{equation}\label{eq:support-translation-puiseux}
  a=c t^q b.
\end{equation}
For an element of the group algebra written uniquely as a finite sum
\[
  u=\sum_{r\in\Q}u_r t^r\in F[\Q],
\]
define its support by
\[
  \operatorname{supp}(u)=\{r\in\Q:u_r\neq0\},
  \qquad \operatorname{supp}(0)=\varnothing.
\]
This definition is used only for group-algebra elements, not for arbitrary
fractions.  With $F[z_n]$ viewed as a subring of $F[\Q]$, both $a$ and $b$
have finite support contained in
$L_n=(1/n!)\Z$.  Multiplication by $t^q$ translates support and causes no
cancellation, so
\[
  \operatorname{supp}(a)=q+\operatorname{supp}(b).
\]
Choose $s\in\operatorname{supp}(b)$.  Then $s\in L_n$ and
$q+s\in L_n$, whence $q\in L_n$.  Therefore $t^q\in z_n^\Z$, so $[f]$ was
already the identity element of $P_n$.  This proves that every $\iota_n$
is injective.

Every rational exponent belongs to some $L_n$, and every element of the
union field belongs to some $F(z_n)$.  Hence the images of the $P_n$ cover
$P(F,t)$.  We may therefore identify
\[
  P(F,t)=\bigcup_{n\geq1}P_n.
\]
For each $n$, choose the explicit free complement $C_n$ constructed above,
so that $P_{n+1}=P_n\oplus C_n$.  Induction gives
\[
  P_m=P_1\oplus\bigoplus_{1\leq n<m}C_n,
\]
and taking the union over $m$ gives
\[
  P(F,t)=P_1\oplus\bigoplus_{n\geq1}C_n.
\]
Both $P_1$ and all the $C_n$ are free abelian; their direct sum is free
abelian.  This proves the lemma.
\end{proof}

\end{document}
